% ============================================================
%  sections/literature_review.tex
%
%  SOURCE MATERIAL: Thesis Chapters 2, 3, and parts of 4
%  JIABR SCOPE: Islamic finance, Shariah governance, fintech,
%               portfolio management — all highly relevant.
%
%  EDITING GUIDE (work section by section with Claude):
%   - Each \subsection can be tackled independently
%   - Add new citations by inserting \citep{key} and adding
%     the entry to references.bib
%   - Keep each paragraph to a clear argumentative point
%   - Target: ~1,800–2,500 words for this section alone
% ============================================================

\section{Literature Review}
\label{sec:lit_review}

\subsection{Islamic Finance: Foundations and Principles}
\label{subsec:islamic_finance}

Islamic finance operates within an ethical framework derived from the Qur'an
and the Sunnah, prohibiting \textit{riba} (interest), \textit{gharar}
(excessive uncertainty), and \textit{maysir} (speculation)
\citep{ayub2009understanding, elgamal2006islamic}. These prohibitions are not
merely procedural constraints; they reflect a broader commitment to what
\citet{asutay2007conceptualisation} terms the Islamic moral economy, wherein
economic activity must promote social welfare and distributive justice. The
industry has grown substantially since the mid-twentieth century, moving from a
niche segment of Middle Eastern banking to a global market with assets
exceeding \$3 trillion \citep{chachi2006thirty}. Despite this growth,
\citet{asutay2012conceptualising} identifies a persistent gap between the
aspirations of Islamic moral economy theory and the actual practice of Islamic
financial institutions, which often mimic conventional products in substance
while maintaining Shariah compliance in form \citep{chong2009islamic}.

Risk-sharing — rather than risk-transfer — is a defining feature that
distinguishes Islamic financial contracts from their conventional counterparts
\citep{chapra2000future}. Products such as \textit{mudaraba} (profit-sharing)
and \textit{musharaka} (joint venture) embody this principle, although their
uptake in practice has been uneven \citep{ahmed2016product}. Within the equity
domain, Shariah screening eliminates firms with impermissible business
activities (alcohol, gambling, conventional finance) or with financial ratios
that exceed prescribed leverage thresholds, creating a restricted investment
universe that may differ substantially in risk-return properties from the
broader market \citep{aysan2016islamic}.

\subsection{Islamic Fintech: Emergence and Applications}
\label{subsec:islamic_fintech}

Financial technology (fintech) — broadly defined as innovation that disrupts or
augments financial services through technology — has reshaped the global
financial landscape since the post-2008 crisis period
\citep{arner2016evolution, arslanian2019future}. Islamic fintech, the
intersection of fintech with Shariah-compliant financial principles, has
emerged as a fast-growing sub-domain with particular relevance for the 1.8
billion Muslims who seek financial products aligned with their religious values
\citep{demirguc2018global}. The potential of Islamic fintech for financial
inclusion is especially pronounced: a significant share of the world's
unbanked adults reside in majority-Muslim countries, where distrust of
interest-bearing products is a documented barrier to formal financial
participation \citep{demirguc2013islamic}.

The applications of Islamic fintech span several domains. Shariah-compliant
digital banking and payment solutions have proliferated, enabling mobile-first
access to profit-and-loss sharing accounts \citep{dawood2022business}.
Crowdfunding and peer-to-peer (P2P) lending platforms, such as Ethis
(\citeyear{ethis2021about}) and Blossom Finance (\citeyear{blossom2023halal}),
have adapted the \textit{musharaka} and \textit{murabaha} structures for
digital distribution. Takatech — the fusion of takaful (Islamic insurance)
with technology — has streamlined mutual insurance products, while robo-advisory
platforms apply algorithmic screening to ensure portfolio compliance with
Shariah standards \citep{cengiz2023fintech}.

Artificial intelligence (AI) and machine learning occupy an increasingly
prominent role within Islamic fintech. \citet{broby2022predictive} documents
the broader shift toward predictive analytics in finance, while
\citet{dawood2022business} identify ML-driven credit scoring and investment
advisory as priority development areas for Islamic fintech firms. The
regulatory dimension is equally important: \citet{ahmad2013regulatory}
highlights how fragmented legal environments across jurisdictions create
compliance burdens for Islamic fintech operators, a challenge that
Regulatory Technology (RegTech) tools — leveraging AI for real-time Shariah
audit — are beginning to address.

\subsection{Open Banking and Its Implications for Islamic Fintech}
\label{subsec:open_banking}

Open banking — the mandated or voluntary sharing of customer financial data
through Application Programming Interfaces (APIs) — has restructured the
competitive landscape of retail banking globally
\citep{briones2022open, borgogno2019data}. The European Union's Payment
Services Directive 2 \citep{eu2015directive} represents the most prominent
regulatory catalyst, while analogous frameworks have emerged in the United
Kingdom, Australia, Brazil \citep{brazil2020open}, and Türkiye
\citep{cbfo2023fintech}. \citet{cbi2023global} documents how open banking
infrastructure reduces entry barriers for fintech challengers by decoupling
financial data from incumbent institutions.

For Islamic fintech, open banking creates both opportunities and risks. On the
opportunity side, API-enabled data sharing supports more granular Shariah
screening, personalised \textit{zakat} (almsgiving) calculations, and
real-time profit-and-loss attribution for equity investment accounts. On the
risk side, data governance frameworks must ensure that shared data is not
employed in ways that breach Islamic ethical standards — for example, in
algorithmic pricing that constitutes \textit{gharar} \citep{cengiz2023fintech}.

\subsection{Machine Learning in Portfolio Optimisation}
\label{subsec:ml_portfolio}

Modern Portfolio Theory (MPT), introduced by \citet{bodie1984nominal} and
formalised in the mean-variance framework of \citeauthor{chen2016xgboost},
remains the canonical approach to portfolio construction. MPT, however, relies
on distributional assumptions (normality of returns, stability of
covariances) that are routinely violated in emerging markets and during periods
of financial stress. Machine learning offers a data-driven alternative capable
of capturing non-linear relationships and temporal dependencies in asset
returns \citep{broby2022predictive}.

Among sequence-based models, LSTM networks — introduced by Hochreiter and
Schmidhuber (1997) and extended to bidirectional architectures — have
demonstrated strong performance in financial time-series forecasting because
their gating mechanisms selectively retain long-range temporal information
\citep{azuaje2006mining}. The Gated Recurrent Unit \citep{cho2014learning}
offers a computationally lighter alternative with comparable forecasting
accuracy. Among non-sequence models, \citet{chen2016xgboost} propose
XGBoost as a scalable gradient-boosted tree system that has become a benchmark
in tabular data competitions. Support Vector Machines \citep{cortes1995support}
have been applied to directional stock prediction with moderate success, while
Random Forest \citep{azuaje2006mining} provides interpretable feature
importance rankings alongside competitive accuracy.

The application of these methods to \textit{Shariah-compliant} equity universes
remains sparse in the published literature. Most existing studies employ
conventional (non-screened) indices, leaving open the question of whether ML
performance advantages generalise to the restricted investment sets imposed by
Islamic screening. This gap is particularly acute for Turkish Islamic equities,
where the BIST-XKTUM Islamic Index provides a well-defined, exchange-regulated
Shariah-compliant benchmark. The present study addresses this gap by
systematically comparing all six ML model families within this specific
investment universe, with portfolio performance evaluated against a classical
Markowitz baseline.

% ── TODO / COLLABORATION NOTES ──────────────────────────────
% ➊ Add any recently published (2023–2025) studies on ML in
%   Islamic finance here — search Scopus for
%   "machine learning" AND "Islamic finance" AND "portfolio".
% ➋ Expand the open banking subsection if Chapter 3 evidence
%   is to be included in the final paper.
% ➌ Review all \citep{} keys against references.bib to ensure
%   no missing entries before compilation.
% ────────────────────────────────────────────────────────────
